\documentclass[a4paper,11pt]{article}

\usepackage{latexsym}
\usepackage{amssymb}
\usepackage{amsthm}
\usepackage{amsmath}
\usepackage{amsrefs}
\usepackage[colorlinks=true,urlcolor=blue,citecolor=blue,linkcolor=blue,linktocpage,pdfpagelabels,bookmarksnumbered,bookmarksopen]{hyperref}
\usepackage{indentfirst}
\usepackage{moreverb}
\usepackage{url}
\pagestyle{myheadings}

\begin{document}


\begin{center}\LARGE\bf
	SMR: LEICOV. Shapes, Materials and Regularity: Lecce Encounters In Calculus Of Variations

\end{center}

\

\begin{center}
	{\bf\large Organizers:}
	Paolo Acampora (Universit\`a del Molise),
	Alessandro Carbotti, Simone Cito, Luigi Negro and Francesco Solombrino (Universit\`a del Salento)
\end{center}

\

\noindent
SMR:LEICOV is a series of three meetings to bring together young researchers and experts in the field of Calculus of Variations and Partial Differential Equations. We want to give value to the intersections between shape optimization, materials science and regularity theory.

\

{\bf List of confirmed speakers (14)}

\begin{enumerate}
	\item Luciana Angiuli, Università del Salento, Italy
	\item Annika Bach, Eindhoven University of Technology, Holland
	\item Dorin Bucur, Université de Savoie, France
	\item Maicol Caponi, Università dell'Aquila, Italy
	\item Filomena De Filippis, University of Salzburg, Austria
	\item Lucia De Luca, CNR, Italy
	\item Manuel Friedrich, Johannes Kepler Universität, Linz, Austria
	\item Ilias Ftouhi, Université de Nîmes, France
	\item Alessandro Giacomini, Università di Brescia, Italy
	\item Ilaria Lucardesi, Università di Pisa, Italy
	\item Carlo Nitsch, Università di Napoli “Federico II”, Italy
	\item Gloria Paoli, Università degli Studi di Napoli "Federico II", Italy
	\item Marianna Porfido, TU Bergakademie Freiberg, Germany
	\item Filippo Riva, Charles University, Prague, Czech Republic
\end{enumerate}

\vspace{0.2 cm}

\textbf{See next pages for titles and abstracts}

\bigskip

\section*{\centering Titles and Abstracts}

\noindent\rule{\textwidth}{0.4pt}
\begin{center}
\subsection*{Hypercontractivity and asymptotic behaviour of generalised Mehler semigroups}

\textbf{Luciana Angiuli} \\
\textit{Università del Salento, Italy} \\
\url{luciana.angiuli@unisalento.it}
\
\end{center}
\medskip
A natural $L^p$-setting for studying semigroups associated with elliptic operators with unbounded coefficients is the one related to invariant measures. This is the case, for instance, of the classical Ornstein–Uhlenbeck semigroup $(P(t))_{t \ge 0}$, which enjoys good properties in $L^p(\gamma)$, where $\gamma$ denotes the standard Gaussian measure that turns out to be the unique invariant measure associated with $(P(t))_{t \ge 0}$. One of the most relevant properties of the Ornstein–Uhlenbeck semigroup, proved by Nelson in 1973, is hypercontractivity: for any $1<p<q<\infty$ there exists $t_0=t_0(p,q)>0$ such that

$$\|P(t)f\|_{L^q(\gamma)} \le \|f\|_{L^p(\gamma)}$$

for all $f \in L^p(\gamma)$ and $t \ge t_0$. The hypercontractivity of $P(t)$ is strictly connected to the validity of the classical logarithmic Sobolev inequality with respect to the Gaussian measure and to the asymptotic behaviour of the semigroup $P(t)$ as $t \to +\infty$.

The Ornstein–Uhlenbeck semigroup is the prototype for the class of generalised Mehler semigroups arising from linear stochastic differential equations perturbed by Lévy noise. For such semigroups, the hypercontractivity property fails to hold in general. In this talk we investigate the hypercontractivity property of generalised Mehler semigroups $(T(t))_{t \ge 0}$ on the $L^p$-scale with respect to an invariant measure $\mu$. When the associated invariant measure $\mu$ lacks a purely Gaussian structure, jump components may prevent the validity of Nelson's classical $L^p$-$L^q$ estimates. However, a summability-improving property can be obtained in the setting of mixed-norm spaces $\mathcal{X}_{p,q}(E;\gamma,\pi)$ related to the factorization of the invariant measure $\mu = \gamma * \pi$ into a Gaussian part $\gamma$ and an infinitely divisible non-Gaussian part $\pi$. As in the classical Gaussian case, some modified logarithmic Sobolev inequalities with respect to invariant measures can be inferred, as well as some results on the asymptotic behaviour of $T(t)$ as $t \to +\infty$.

\smallskip
Joint work with Simone Ferrari and Diego Pallara.

\newpage

\noindent\rule{\textwidth}{0.4pt}

\begin{center}
\subsection*{Monotonicity arguments for the multi-cell formula in stochastic homogenization}

\textbf{Annika Bach} \\
\textit{Eindhoven University of Technology, Holland} \\
\url{a.bach@tue.nl}
\
\end{center}
\medskip
Stochastic homogenization of integral functionals $F_\varepsilon$ via $\Gamma$-convergence can be seen as a rigorous mathematical averaging procedure, where highly oscillating random integrands are in the $\Gamma$-limit replaced by a homogeneous (suitably averaged) integrand. This limiting integrand is typically determined via a multi-cell formula, a limit of minimization problems where the integral functionals $F_\varepsilon$ are rescaled to larger and larger cubes and subjected to suitable boundary conditions. In many situations it is useful to replace those minimization problems by a family of auxiliary problems defined on a restricted but growing class of competitors, and to relate the two minimization problems via a monotonicity argument. In this talk we will present two such monotonicity arguments that we used to obtain fluctuation estimates for the multi-cell formula in the context of integral functionals defined on finite partitions and to obtain a stochastic homogenization result for integral functionals defined on manifold-valued Sobolev functions.

\smallskip
Joint work with Matthias Ruf (Justus-Liebig-Universität Gießen).

\newpage

\noindent\rule{\textwidth}{0.4pt}

\begin{center}
\subsection*{On a free boundary problem for the gradient magnitude of the torsion function}

\textbf{Dorin Bucur} \\
\textit{Université de Savoie, France} \\
\url{dorin.bucur@univ-savoie.fr}
\
\end{center}
\medskip
A recent result by Burdzy, Ftouhi, and Mariano establishes a series of geometric properties of optimizers among convex sets for the maximum of the gradient of the torsion function. In this talk, I will discuss a related free boundary problem. We prove a general result concerning the presence of flat parts in the free boundary, together with an estimate of their Hausdorff measure. We will also discuss applications to shape optimization problems.

\newpage

\noindent\rule{\textwidth}{0.4pt}

\begin{center}
\subsection*{$H$-convergence and $\Gamma$-convergence in the Riesz fractional setting: the nonlinear case}

\textbf{Maicol Caponi} \\
\textit{Università dell'Aquila, Italy} \\
\url{maicol.caponi@univaq.it}
\
\end{center}
\medskip
This talk concerns the $H$-convergence of nonlinear nonlocal monotone operators defined through the Riesz fractional gradient and divergence. We show that the $H$-convergence in this nonlocal framework is equivalent to the $H$-convergence of the corresponding local one. As a consequence, we obtain a $H$-compactness result for a suitable class of nonlocal monotone operators. We then study the $\Gamma$-convergence of nonlocal energy functionals associated with the subclass of conservative monotone operators, proving that it is equivalent to the $\Gamma$-convergence of the corresponding local energies. A key ingredient is a new uniqueness result for the integral representation of both local and nonlocal functionals. As a by-product, we obtain the $\Gamma$-compactness of the class of nonlocal energies under consideration. Finally, we show the equivalence between the $H$-convergence of nonlocal conservative monotone operators and the $\Gamma$-convergence of the associated energy functionals.

\smallskip
Joint work with G. C. Brusca (SISSA), A. Carbotti (Università del Salento), A. Maione (Politecnico di Milano) and F. Paronetto (Università di Padova).

\newpage

\noindent\rule{\textwidth}{0.4pt}

\begin{center}
\subsection*{Intrinsic Schauder estimates at nearly linear growth}

\textbf{Filomena De Filippis} \\
\textit{University of Salzburg, Austria} \\
\url{filomena.defilippis@plus.ac.at}
\
\end{center}
\medskip
Variational integrals at nearly linear growth appear in the theory of plasticity with logarithmic hardening, that is the borderline configuration between plasticity with power hardening and perfect plasticity. The related (very challenging) regularity theory for minima has been intensively developed over the last 25 years, see e.g. the work of Frehse \& Seregin '99, Fuchs \& Mingione '00, Bildhauer '03, Beck \& Bulíček \& Gmeineder '20, Di Marco \& Marcellini '20, Gmeineder \& Kristensen '22, De Filippis \& Mingione '23.

In this talk, we will discuss a nonlinear potential theoretic framework for Schauder estimates for vector valued solutions of a broad class of nonautonomous variational problems at nearly linear growth. Our approach naturally embraces the variable exponent as well as the Double and Multi phase setting, yielding new regularity results in basic models and recovering optimal regularity recently established in specific cases.

\smallskip
Joint work with Cristiana De Filippis (Parma) and Peter Hästö (Helsinki).

\newpage

\noindent\rule{\textwidth}{0.4pt}

\begin{center}
\subsection*{Asymptotic expansion of the relative $s$-fractional perimeter as $s$ tends to 0}

\textbf{Lucia De Luca} \\
\textit{CNR, Italy} \\
\url{l.deluca@iac.cnr.it}
\
\end{center}
\medskip
We prove a first-order $\Gamma$-convergence result (as $s$ tends to 0) for the $s$-fractional perimeter relative to a given bounded set and with prescribed exterior datum. We investigate the robustness of our approach under perturbations of the exterior datum and in presence of volume constraints. We then use the $\Gamma$-convergence expansion to study the asymptotic behavior of $s$-minimal sets as $s$ tends to 0.

\smallskip
Joint work with G. Palatucci, M. Piccinini and M. Ponsiglione.

\newpage

\noindent\rule{\textwidth}{0.4pt}

\begin{center}
\subsection*{Quasistatic evolution of cohesive-type fracture}

\textbf{Manuel Friedrich} \\
\textit{Johannes Kepler Universität, Linz, Austria} \\
\url{manuel.friedrich@jku.at}
\
\end{center}
\medskip
In this talk, I present an existence result for a variational model of quasi-static crack growth for a class of cohesive fracture energies with an activation threshold. The evolution is characterized by an irreversibility condition in terms of a maximal opening variable, a global stability, and an energy balance.

\smallskip
Joint work with Vito Crismale.

\newpage

\noindent\rule{\textwidth}{0.4pt}

\begin{center}
\subsection*{About the maximization of the gradient of the torsion function}

\textbf{Ilias Ftouhi} \\
\textit{Université de Nîmes, France} \\
\url{ilias.ftouhi@unimes.fr}
\
\end{center}
\medskip
In this talk, we address the problem of maximizing the supremum norm of the gradient of the torsion function over planar convex domains with fixed measure (or perimeter). We prove that an optimal shape exists, its boundary possesses $C^1$ regularity, and it features at least one straight segment. The proofs leverage probabilistic methods along with a novel formulation of the boundary Harnack principle for the torsion function.

\smallskip
Joint work with Krzysztof Burdzy (University of Washington) and Phanuel Mariano (Union College).

\newpage

\noindent\rule{\textwidth}{0.4pt}

\begin{center}
\subsection*{Periodic homogenization in scalar Hencky plasticity}

\textbf{Alessandro Giacomini} \\
\textit{Università di Brescia, Italy} \\
\url{alessandro.giacomini@unibs.it}
\
\end{center}
\medskip
In this talk I will deal with periodic homogenization for a two-phase elasto-plastic material in the framework of $\Gamma$-convergence under two-scale convergence. The elasto-plastic character of the homogenized problem is unclear, and will be discussed by resorting to a one field description of the associated energy. A positive answer will be given in one dimension.

\smallskip
Joint work with G. Francfort.

\newpage

\noindent\rule{\textwidth}{0.4pt}

\begin{center}
\subsection*{On the shape optimization of the torsional rigidity under thickness constraint among triangles}

\textbf{Ilaria Lucardesi} \\
\textit{Università di Pisa, Italy} \\
\url{ilaria.lucardesi@unipi.it}
\
\end{center}
\medskip
In this talk I will present some recent results about the optimization, both minimization and maximization, of a shape functional which depends on the torsional rigidity and the thickness, in the class of triangles. Besides the detection of optimal triangles, the study allows to obtain new geometric estimates on the torsional rigidity for triangles.

\smallskip
Joint work with Davide Zucco (Università di Torino).

\newpage

\noindent\rule{\textwidth}{0.4pt}

\begin{center}
\subsection*{Improving a Spectral Inequality by Payne}

\textbf{Carlo Nitsch} \\
\textit{Università di Napoli “Federico II”, Italy} \\
\url{c.nitsch@unina.it}
\
\end{center}
\medskip
A celebrated inequality by Payne relates the first eigenvalue of the Dirichlet Laplacian to the first eigenvalue of the buckling problem. Motivated by the goal of establishing a quantitative version of this inequality, we show that Payne's original estimate — which is not sharp — can in fact be improved. Our result provides a refined spectral bound and opens the way to further investigations into quantitative enhancements of classical inequalities in spectral theory.

\newpage

\noindent\rule{\textwidth}{0.4pt}

\begin{center}
\subsection*{Title to be announced}

\textbf{Gloria Paoli} \\
\textit{Università degli Studi di Napoli "Federico II", Italy} \\
\url{gloria.paoli@unina.it}
\
\end{center}
\medskip
Abstract to be announced.

\newpage

\noindent\rule{\textwidth}{0.4pt}

\begin{center}
\subsection*{The Free Lunch Theorem of homogenisation}

\textbf{Marianna Porfido} \\
\textit{TU Bergakademie Freiberg, Germany} \\
\url{Marianna.Porfido@math.tu-freiberg.de}
\
\end{center}
\medskip
In this talk, we revisit the classical elliptic homogenisation theory for local coefficients, i.e., $L_\infty$-multiplication operators (H-convergence as envisioned by Murat and Tartar in the 1970's). Then, we compare it with an abstract operator-theory approach that allows for nonlocal coefficients, i.e., general bounded operators (nonlocal H-convergence introduced in [1]). With these concepts in hand, we show that H-convergence for multiplication type operators always implies nonlocal H-convergence. As a consequence, H-convergence for multiplication operators in divergence form problems will always imply H-type convergence for a different variational problem for free.

\smallskip
Joint work with Andreas Buchinger and Marcus Waurick.

\medskip

\begin{center} {\bf References} \end{center}

\noindent{[1] M. Waurick, Nonlocal H-Convergence, Calc. Var. Partial Differential Equations 57(6) (2018).}\\
\noindent{[2] A. Buchinger, M. Porfido, M. Waurick, The Free Lunch Theorem of Homogenisation, preprint (2026), doi: https://arxiv.org/abs/2606.04498}

\newpage

\noindent\rule{\textwidth}{0.4pt}

\begin{center}
\subsection*{A free boundary approach to quasistatic debonding}

\textbf{Filippo Riva} \\
\textit{Charles University, Prague, Czech Republic} \\
\url{filippo.riva@matfyz.cuni.cz}
\
\end{center}
\medskip
Debonding models describe the evolution of an adhesive elastic membrane peeled away from a planar rigid substrate. When the process is sufficiently slow, inertial and viscous effects can be neglected, leading to a quasistatic setting. In this framework, debonding can be formulated as a variational evolution of sets. In the talk, I will present how to prove existence of so-called energetic solutions for this formulation — based on global minimizers of a suitable functional together with an energy balance — by exploiting an equivalent rewriting of the model in terms of the celebrated one-phase Bernoulli free boundary problem.

\smallskip
Joint work with E. Maggiorelli and E. Tolotti.

\end{document}
